

\medskip
\noindent\SGAInlineRefTarget{sga3:ii:proposition:4-5}\textbf{Proposition 4.5.}
\textit{Let \(F\) be a good \(\mathcal O_S\)-module. Then:}

\textit{(i) The \(S\)-group
\(\underline{\operatorname{Aut}}_{\mathcal O_S\text{-mod.}}(F)\)
satisfies condition \((E)\), and there is a functorial isomorphism}
\[
 \underline{\operatorname{Lie}}\!
   \left(\underline{\operatorname{Aut}}_{\mathcal O_S\text{-mod.}}(F)/S,
         \mathcal M\right)
 \xrightarrow{\sim}
 \underline{\operatorname{Hom}}_{\mathcal O_S\text{-mod.}}
   \left(F,\underline{\operatorname{Lie}}(F/S,\mathcal M)\right),
\]
\textit{which respects the \(\mathcal O_S\)-module structures obtained from
condition \((E)\). In particular, there is an isomorphism of
\(\mathcal O_S\)-modules}
\[
 \underline{\operatorname{Lie}}\!
   \left(\underline{\operatorname{Aut}}_{\mathcal O_S\text{-mod.}}(F)/S\right)
 \xrightarrow{\sim}
 \underline{\operatorname{End}}_{\mathcal O_S\text{-mod.}}(F).
\]

\textit{\SGAInlineRefTarget{sga3:ii:proposition:4-5:item:ii}(ii)\footnote{Editor's note: part \SGARefLink{sga3:ii:proposition:4-5:item:ii}{(ii)}, an immediate consequence
of what precedes, has been added.} Moreover,
\(\underline{\operatorname{End}}_{\mathcal O_S\text{-mod.}}(F)\)
is a good \(\mathcal O_S\)-module.}

Indeed, by \SGARefLink{sga3:ii:corollary:4-4-1}{4.4.1}, \(F/S\) satisfies \((E)\), and
\begin{equation*}
 \underline{\operatorname{Lie}}(F/S,\mathcal M)
 =\underline{\operatorname{Lie}}'(F/S,\mathcal M)
 \simeq
 F\otimes_{\mathcal O_S}
 \underline{\operatorname{Lie}}(\mathbf O_S/S,\mathcal M).
 \tag{1}\label{sga3:ii:proposition:4-5:equation:1}
\end{equation*}
Assertion (i) now follows from \SGARefLink{sga3:ii:corollary:4-3-1}{4.3.1}. Put
\(E=\underline{\operatorname{End}}_{\mathcal O_S\text{-mod.}}(F)\).
By (1) and \SGARefLink{sga3:ii:remark:4-3-5}{Remark 4.3.5}, one has the following commutative diagram of
morphisms of abelian \(S\)-groups.

\SGAThreeNativeDiagram[0.90\textwidth]
  {native_diagrams/exp2/exp2_p034_endomorphism_differential_square.tex}
  {The comparison diagram for the differential \(d_E\).}
  {fig:exp2-p45-endomorphism-differential}
  {Polo--Gille re-edition, Expos\'e II}

Here \(d_F\) and the vertical isomorphism \((*)\) are isomorphisms of
\(\mathcal O_S\)-modules; consequently, so is \(d_E\). This proves \SGARefLink{sga3:ii:proposition:4-5:item:ii}{(ii)}.

\medskip
\noindent\SGAInlineRefTarget{sga3:ii:scholium:4-5-1}\textbf{Scholium 4.5.1.}
\footnote{Editor's note: this scholium, implicit in the original, has been
added for use in \SGARefLink{sga3:ii:subsection:4-7-1}{4.7.1}. Here and below, \(t,t'\), and so forth denote
square-zero variables; the letter \(\varepsilon\) is reserved for the unit
section of a group.}
Put (cf. \SGARefLink{sga3:ii:definition:2-1}{2.1})
\(\mathcal O_{I_S}=\mathcal O_S\oplus t\mathcal O_S\), with \(t^2=0\),
and let \(F\) be a good \(\mathcal O_S\)-module. For every \(S'\to S\),
the morphism
\[
 F(S')\oplus tF(S')
 =F(S')\otimes_{\mathcal O(S')}\mathcal O(I_{S'})
 \longrightarrow
 F(I_{S'})
 =F(S')\oplus\underline{\operatorname{Lie}}(F/S)(S'),
\]
which is the identity on \(F(S')\), induces an isomorphism of
\(\mathcal O(S')\)-modules
\[
 tF(S')\xrightarrow{\sim}\underline{\operatorname{Lie}}(F/S)(S').
\]
As \(S'\) varies, this gives an isomorphism that we may write
\[
 \underline{\operatorname{Lie}}(F/S)\simeq tF.
\]

Thus, for every \(S'\to S\), \SGARefLink{sga3:ii:proposition:4-5}{Proposition 4.5} gives a commutative
diagram.

\SGAThreeNativeDiagram[0.91\textwidth]
  {native_diagrams/exp2/exp2_p035_dual_number_automorphism_square.tex}
  {Endomorphisms and automorphisms over the dual numbers.}
  {fig:exp2-p451-dual-number-automorphisms}
  {LNM 151, pp.~66--67}

It follows from 4.3 bis that every
\[
 X\in
 \underline{\operatorname{End}}_{\mathcal O_{S'}\text{-mod.}}(F_{S'})
\]
corresponds to
\(\mathrm{id}+tX\in
\underline{\operatorname{Aut}}_{\mathcal O_{I_{S'}}\text{-mod.}}
(F_{I_{S'}})\).

\medskip
\noindent\SGAInlineRefTarget{sga3:ii:definition:4-6}\textbf{Definition 4.6.}
An \(S\)-functor in groups \(G\) is called \emph{good} if \(G/S\)
satisfies condition \((E)\) and
\(\underline{\operatorname{Lie}}(G/S)\) is a good
\(\mathcal O_S\)-module.

\footnote{Editor's note: the following sentence has been added.}
Notice that if \(F\) is a good \(\mathcal O_S\)-module, then it is a
good \(S\)-group: indeed, \(F/S\) satisfies \((E)\), and
\(\underline{\operatorname{Lie}}(F/S)\simeq F\) is a good
\(\mathcal O_S\)-module.

\medskip
\noindent\SGAInlineRefTarget{sga3:ii:example:4-6-1}\textbf{Example 4.6.1.}
\textit{If \(G\) is representable, it is good. Indeed, \(G/S\)
satisfies \((E)\), and \(\underline{\operatorname{Lie}}(G/S)\) has
the form \(V(E)\), hence is good by \SGARefLink{sga3:ii:examples:4-4-2}{4.4.2}.}

\medskip
\noindent\SGAInlineRefTarget{sga3:ii:lemma:4-6-2}\textbf{Lemma 4.6.2.}
\footnote{Editor's note: this lemma has been added.}
\textit{Let \(G\) be an \(S\)-functor in groups such that \(G/S\)
satisfies \((E)\), and put
\(F=\underline{\operatorname{Lie}}(G/S)\). Then \(F\) satisfies
\((E)\), and the morphism of abelian groups}
\[
 d:F\longrightarrow\underline{\operatorname{Lie}}(F/S)
\]
\textit{respects the \(\mathcal O_S\)-module structures. Consequently,
\(G\) is good if and only if
\(F\to\underline{\operatorname{Lie}}(F/S)\) is bijective.}

\textit{Proof.}
Suppose that \(G/S\) satisfies \((E)\). Let
\(\mathcal M,\mathcal N\) be finite free \(\mathcal O_S\)-modules,
write
\[
 F(\mathcal N)=\underline{\operatorname{Lie}}(G/S,\mathcal N),
\]
and let \(e\) denote the unit section of \(G\). For every \(S'\to S\),
\[
 F(\mathcal N)(S')
 =
 \left\{
  g\in\operatorname{Hom}_S(I_{S'}(\mathcal N),G)
  \ \middle|\ g\circ\varepsilon_{\mathcal N}=e
 \right\}.
\]
Moreover,
\(\underline{\operatorname{Lie}}(F(\mathcal N)/S,\mathcal M)(S')\)
identifies with the set of morphisms
\[
\left\{
\begin{array}{@{}l@{\ }|@{\ }l@{}}
 \Phi\in
 \operatorname{Hom}_S
   (I_{S'}(\mathcal N)\times_{S'}I_{S'}(\mathcal M),G)
 &
 \begin{aligned}
 \Phi\circ
  (\varepsilon_{\mathcal N}\times
   \mathrm{id}_{I_{S'}(\mathcal M)})
   &=e\in G(I_{S'}(\mathcal M)),\\
 \Phi\circ
  (\mathrm{id}_{I_{S'}(\mathcal N)}\times
   \varepsilon_{\mathcal M})
   &=e\in G(I_{S'}(\mathcal N)).
 \end{aligned}
\end{array}
\right\}.
\]
This proves
\begin{equation*}
 \underline{\operatorname{Lie}}(F(\mathcal N)/S,\mathcal M)
 \simeq
 \underline{\operatorname{Lie}}(F(\mathcal M)/S,\mathcal N).
 \tag{1}\label{sga3:ii:lemma:4-6-2:equation:1}
\end{equation*}
Since \(G/S\) satisfies \((E)\), it follows that
\[
\begin{aligned}
 \underline{\operatorname{Lie}}
   (F(\mathcal N)/S,\mathcal M_1\oplus\mathcal M_2)
 &\simeq
 \underline{\operatorname{Lie}}
   (F(\mathcal M_1\oplus\mathcal M_2)/S,\mathcal N)\\
 &\simeq
 \underline{\operatorname{Lie}}
   ((F(\mathcal M_1)\times_SF(\mathcal M_2))/S,\mathcal N).
\end{aligned}
\]
By \SGARefLink{sga3:ii:proposition:3-8}{3.8}, the last term is isomorphic to
\[
\begin{aligned}
 &\underline{\operatorname{Lie}}(F(\mathcal M_1)/S,\mathcal N)
 \times_S
 \underline{\operatorname{Lie}}(F(\mathcal M_2)/S,\mathcal N)\\
 &\qquad\simeq
 \underline{\operatorname{Lie}}(F(\mathcal N)/S,\mathcal M_1)
 \times_S
 \underline{\operatorname{Lie}}(F(\mathcal N)/S,\mathcal M_2).
\end{aligned}
\]

Hence
\begin{equation*}
 \underline{\operatorname{Lie}}
   (F(\mathcal N)/S,\mathcal M_1\oplus\mathcal M_2)
 \simeq
 \underline{\operatorname{Lie}}(F(\mathcal N)/S,\mathcal M_1)
 \times_S
 \underline{\operatorname{Lie}}(F(\mathcal N)/S,\mathcal M_2),
 \tag{2}\label{sga3:ii:lemma:4-6-2:equation:2}
\end{equation*}
so \(F(\mathcal N)/S\) satisfies \((E)\), by Remark 4.1.1.2(a).

It remains to show that the morphism of abelian groups
\[
 d:F(\mathcal N)\longrightarrow
 \underline{\operatorname{Lie}}(F(\mathcal N)/S)
\]
respects the \(\mathcal O_S\)-module structures. Consider the free
\(\mathcal O_S\)-module \(\mathcal M=t\mathcal O_S\). Then
\[
 \mathcal O(I_{S'}(\mathcal M))
 =\mathcal O(S')[t]/(t^2)
 =\mathcal O(S')\oplus t\mathcal O(S'),
 \qquad
 \underline{\operatorname{Lie}}(\mathbf O_S/S)\simeq t\mathcal O_S.
\]
Let \(\rho_t:I_S(\mathcal M)\to S\) be the structural morphism,
corresponding to the injection
\[
 u_t:\mathcal O_S\hookrightarrow
 D_{\mathcal O_S}(\mathcal M)
 =\mathcal O_S\oplus t\mathcal O_S.
\]
Recall (cf. 3.4.2) that for every \(S\)-functor \(X\),
\(F(\mathcal N)(X)\) is the set of \(S\)-morphisms
\(\phi:I_X(\mathcal N)\to G\) such that
\(\phi\circ\varepsilon_{\mathcal N}=e\), and the action of
\(a\in\mathcal O(X)\) is
\[
 a\cdot\phi=\phi\circ a^*,
\]
where \(a^*\) is the endomorphism of \(I_X(\mathcal N)\) associated
with \(a\); see \SGARefLink{sga3:ii:subsection:2-1-3}{2.1.3}.

Consequently, by \SGARefLink{sga3:ii:remark:4-3-4}{4.3.4}, the composite
\[
 d:F(\mathcal N)\xrightarrow{\sim}
 F(\mathcal N)\otimes_{\mathcal O_S}t\mathcal O_S
 \longrightarrow
 \underline{\operatorname{Lie}}(F(\mathcal N)/S)
\]
is given, for \(S'\to S\) and
\(f\in\operatorname{Hom}_S(S',F(\mathcal N))\), by
\[
 f\longmapsto f\otimes t
 \longmapsto t\cdot(f\circ\rho_t)
 =f\circ\rho_t\circ t^*.
\]
Here \(f\circ\rho_t\in F(\mathcal N)(I_{S'}(\mathcal M))\) is viewed
as an \(S\)-morphism
\[
 \phi:I_{I_{S'}(\mathcal M)}(\mathcal N)\longrightarrow G
\]
such that \(\phi\circ\varepsilon_{\mathcal N}=e\). We must prove that
\(d\) is \(\mathcal O_S\)-linear, or equivalently that
\[
 d(a\cdot f)=u_t(a)\cdot d(f)
 \qquad(a\in\mathcal O(S')).
\]
Viewing \(f\in F(\mathcal N)(S')\) as a morphism
\(I_{S'}(\mathcal N)\to G\), one likewise has
\(a\cdot f=f\circ a^*\). Functoriality in \(X\) of the action of
\(\mathcal O(X)\) on \(I_X(\mathcal N)\) gives the following
commutative diagram.

\SGAThreeNativeDiagram[0.57\textwidth]
  {native_diagrams/exp2/exp2_p036_scalar_action_square.tex}
  {Compatibility of scalar action with the structural morphism \(\rho_t\).}
  {fig:exp2-p462-scalar-action}
  {Polo--Gille re-edition, Expos\'e II}

Therefore
\[
\begin{aligned}
 d(a\cdot f)
 &=f\circ a^*\circ\rho_t\circ t^*\\
 &=f\circ\rho_t\circ u_t(a)^*\circ t^*\\
 &=f\circ\rho_t\circ t^*\circ u_t(a)^*
  =u_t(a)\cdot d(f).
\end{aligned}
\]
The penultimate equality follows from the commutativity of
\(\mathcal O\). This completes the proof of \SGARefLink{sga3:ii:lemma:4-6-2}{Lemma 4.6.2}.

\medskip
\noindent\SGAInlineRefTarget{sga3:ii:theorem:4-7}\textbf{Theorem 4.7.}
\textit{If \(F\) is a good \(\mathcal O_S\)-module, then the
\(S\)-group
\(\underline{\operatorname{Aut}}_{\mathcal O_S\text{-mod.}}(F)\)
is good.}

\footnote{Editor's note: the original has been simplified here by using
the additions made in \SGARefLink{sga3:ii:corollary:4-3-1}{4.3.1} and \SGARefLink{sga3:ii:proposition:4-5}{4.5}.}
Indeed, by \SGARefLink{sga3:ii:proposition:4-5}{4.5},
\(\underline{\operatorname{Aut}}_{\mathcal O_S\text{-mod.}}(F)/S\)
satisfies \((E)\), while
\[
 \underline{\operatorname{Lie}}\!
  \left(\underline{\operatorname{Aut}}_{\mathcal O_S\text{-mod.}}(F)/S\right)
 \simeq
 \underline{\operatorname{End}}_{\mathcal O_S\text{-mod.}}(F)
\]
is a good \(\mathcal O_S\)-module.

\subsection*{\SGAInlineRefTarget{sga3:ii:subsection:4-7-1}4.7.1.}
Let \(G\) be an \(S\)-group and \(F\) a good
\(\mathcal O_S\)-module. Suppose given a linear representation of
\(G\) in \(F\), that is (\SGARefLink{sga3:i:definition:4-7-1}{I~4.7.1}), a morphism of \(S\)-groups
\[
 \rho:G\longrightarrow
 \underline{\operatorname{Aut}}_{\mathcal O_S\text{-mod.}}(F).
\]
If \(G/S\) satisfies \((E)\), then \SGARefLink{sga3:ii:remark:4-1-c}{4.1.C} and \SGARefLink{sga3:ii:proposition:4-5}{4.5} give a morphism of
\(\mathcal O_S\)-modules, denoted \(\rho'\) or \(d\rho\),
\[
 \underline{\operatorname{Lie}}(G/S)
 \longrightarrow
 \underline{\operatorname{Lie}}\!
  \left(\underline{\operatorname{Aut}}_{\mathcal O_S\text{-mod.}}(F)/S\right)
 \simeq
 \underline{\operatorname{End}}_{\mathcal O_S\text{-mod.}}(F).
\]
\footnote{Editor's note: the following sentence has been added.}
Furthermore, put
\(\mathcal O_{I_S}=\mathcal O_S\oplus t\mathcal O_S\), with
\(t^2=0\). It follows from \SGARefLink{sga3:ii:scholium:4-5-1}{4.5.1} that, if \(S'\to S\) and
\[
 X\in\underline{\operatorname{Lie}}(G/S)(S')\subset G(I_{S'}),
\]
then in
\(\underline{\operatorname{Aut}}_{\mathcal O_{I_{S'}}\text{-mod.}}
(F_{I_{S'}})\) one has
\begin{equation*}
 \rho(X)=\mathrm{id}+t\rho'(X).
 \tag{*}\label{sga3:ii:subsection:4-7-1:equation:star}
\end{equation*}
In other words, for every \(S''\to I_{S'}\) and \(f\in F(S'')\),
\[
 \rho(X)(f)=f+t\rho'(X)(f)
 \quad\text{in }F(S'').
\]

\medskip
\noindent\SGAInlineRefTarget{sga3:ii:definition:4-7-2}\textbf{Definition 4.7.2.}
Let \(G\) be a good \(S\)-group. Then
\(\underline{\operatorname{Lie}}(G/S)\) is a good
\(\mathcal O_S\)-module, and there is a morphism of \(S\)-groups
\[
 \operatorname{Ad}:G\longrightarrow
 \underline{\operatorname{Aut}}_{\mathcal O_S\text{-mod.}}
   \bigl(\underline{\operatorname{Lie}}(G/S)\bigr).
\]
By \SGARefLink{sga3:ii:subsection:4-7-1}{4.7.1}, it induces a morphism of \(\mathcal O_S\)-modules
\[
 \operatorname{ad}:
 \underline{\operatorname{Lie}}(G/S)
 \longrightarrow
 \underline{\operatorname{End}}_{\mathcal O_S\text{-mod.}}
   \bigl(\underline{\operatorname{Lie}}(G/S)\bigr),
\]
or, equivalently, an \(\mathcal O_S\)-bilinear morphism
\[
 \underline{\operatorname{Lie}}(G/S)
 \times_S
 \underline{\operatorname{Lie}}(G/S)
 \longrightarrow
 \underline{\operatorname{Lie}}(G/S),
 \qquad
 (x,y)\longmapsto[x,y]=\operatorname{ad}(x)\cdot y.
\]
Here \(x,y\) are arbitrary elements of
\[
 \underline{\operatorname{Lie}}(G/S)(S')
 =
 \underline{\operatorname{Lie}}(G_{S'}/S')(S').
\]
If \(G\) is commutative, then \([x,y]=0\).

\subsection*{\SGAInlineRefTarget{sga3:ii:subsection:4-7-3}4.7.3.}
The bracket admits the following equivalent description. It is enough
to consider sections over \(S\), that is,
\[
 x,y\in\underline{\operatorname{Lie}}(G/S)(S).
\]
There is a canonical isomorphism
\(I_S\times_SI_S\simeq I_{I_S}\). To avoid confusion, let \(I\) and
\(I'\) denote two copies of \(I_S\), and put
\[
 \mathcal O_I=\mathcal O_S[t],
 \qquad
 \mathcal O_{I'}=\mathcal O_S[t'],
 \qquad
 t^2=0=(t')^2.
\]
There is then a commutative square

\SGAThreeNativeDiagram[0.32\textwidth]
  {native_diagrams/exp2/exp2_p037_dual_number_base_square.tex}
  {The two dual-number structures on \(I\times_S I'\).}
  {fig:exp2-p473-dual-number-base}
  {LNM 151, pp.~68--69}

in which the two arrows from \(I\times I'\) identify it with the
scheme of dual numbers over \(I\) or over \(I'\), respectively.
Writing \(L=\underline{\operatorname{Lie}}(G/S)\), one obtains the
following commutative diagram of groups.

\SGAThreeNativeDiagram[0.69\textwidth]
  {native_diagrams/exp2/exp2_p038_exact_sequence_diagram_1.tex}
  {The first exact-sequence diagram used to construct the bracket.}
  {fig:exp2-p473-exact-sequences-one}
  {LNM 151, pp.~69--70}

The ninth object in this diagram is
\(\underline{\operatorname{Lie}}(L/S)(S)\). If \(G\) is good, this
is \(L(S)\), and hence one obtains the next commutative diagram. Its
rows and columns are exact sequences of groups; the five groups
denoted \(L(-)\) are commutative.\footnote{Editor's note: the original
has been expanded in what follows.} Under the identifications
\[
 L(I)=L(S)\oplus tL(S),
 \qquad
 L(I')=L(S)\oplus t'L(S),
\]
the injections \(L(S)\hookrightarrow L(I)\) and
\(L(S)\hookrightarrow L(I')\) are \(u\mapsto tu\) and
\(u\mapsto t'u\), respectively.

\SGAThreeNativeDiagram[0.69\textwidth]
  {native_diagrams/exp2/exp2_p038_exact_sequence_diagram_2.tex}
  {The exact-sequence diagram defining \([x,y]\).}
  {fig:exp2-p473-exact-sequences-two}
  {LNM 151, pp.~69--70}

In such a diagram, take \(x,y\) as indicated and choose arbitrary
lifts \(\widetilde x\in L(I)\) of \(x\) and \(y'\in L(I')\) of \(y\).
The commutator
\[
 \widetilde x\,y'\,\widetilde x^{-1}(y')^{-1}
 \quad\text{in }G(I\times I')
\]
is independent of the chosen lifts and is the image of a uniquely
determined element \(z\) as shown. The reader will verify that \(z=[x,y]\).\footnote{The writer
acknowledges, at Gabriel's request, that the exercise is not immediate;
this is, moreover, why it is not given in the text.}
\footnote{Editor's note: the original said in a footnote that the author,
at Gabriel's request, acknowledged that the exercise was not immediate,
which was precisely why it had not been placed in the text. The details
have been supplied here, taking account of the addition in \SGARefLink{sga3:ii:subsection:4-7-1}{4.7.1}; see
also [DG70], \(\S\)II.4, \SGARefLink{sga3:ii:subsection:4-2}{4.2}.}

Indeed, use the same letters \(x\) and \(y\) for their images under
the canonical sections \(L(S)\to L(I)\) and \(L(S)\to L(I')\).
Then
\[
 \widetilde x=xu,\qquad y'=yv,
 \qquad u,v\in L(S)=L(I)\cap L(I').
\]
Since \(L(I)\) and \(L(I')\) are commutative,
\[
\begin{aligned}
 \widetilde x\,y'\,\widetilde x^{-1}(y')^{-1}
 &=xuyvu^{-1}x^{-1}v^{-1}y^{-1}\\
 &=xuyu^{-1}vx^{-1}v^{-1}y^{-1}\\
 &=xyx^{-1}y^{-1}.
\end{aligned}
\]

This element maps to the unit in both \(G(I)\) and \(G(I')\), and
therefore comes from the unique \(z\) indicated above. Finally,
viewing \(y\) as an element of \(L_{I'}(I')\), and \(x\) as an
element of \(L(S)\subset G(I')\), \SGARefLink{sga3:ii:subsection:4-7-1}{formula 4.7.1} \((*)\) gives
\[
 xyx^{-1}
 =\operatorname{Ad}(x)(y)
 =(\mathrm{id}+t'\operatorname{ad}(x))(y)
 =y+t'[x,y].
\]
Thus \(xyx^{-1}y^{-1}\in L(I')\) is the image of
\(z=[x,y]\in L(S)\).

Two properties are now apparent from this construction:

\textit{\SGAInlineRefTarget{sga3:ii:subsection:4-7-3:item:i}(i)} the bracket is functorial in \(G\). More precisely,
\[
 G\longmapsto\underline{\operatorname{Lie}}(G/S)
\]
is a functor from good \(S\)-groups to good \(\mathcal O_S\)-modules
equipped with an \(\mathcal O_S\)-bilinear composition law;

\textit{(ii)} \([x,y]+[y,x]=0\), since the diagram is symmetric about
its main diagonal.\footnote{Editor's note: that is, \(x\) and \(y\)
play symmetric roles. The image of \([y,x]\) in \(G(I\times I')\) is
the commutator
\(y'\widetilde x(y')^{-1}\widetilde x^{-1}\), the inverse of
\(\widetilde x y'\widetilde x^{-1}(y')^{-1}=[x,y]\). On the other
hand (cf. \SGARefLink{sga3:ii:definition:4-10}{4.10} below), it is not known at this point whether
\([x,x]=0\) necessarily holds: if \(x\) is represented by
\(\widetilde x\in L(I)\) and by \(x'\in L(I')\), it is not a priori
clear that \(\widetilde x\) and \(x'\) commute.}

\medskip
\noindent\SGAInlineRefTarget{sga3:ii:proposition:4-8}\textbf{Proposition 4.8.}
\textit{Let \(F\) be a good \(\mathcal O_S\)-module. Under the
identification}
\[
 \underline{\operatorname{Lie}}\!
  \left(\underline{\operatorname{Aut}}_{\mathcal O_S\text{-mod.}}(F)/S\right)
 =
 \underline{\operatorname{End}}_{\mathcal O_S\text{-mod.}}(F),
\]
\textit{one has}
\[
 \operatorname{Ad}(g)\cdot Y=g\circ Y\circ g^{-1},
 \qquad
 [X,Y]=X\circ Y-Y\circ X,
\]
\textit{for every \(S'\to S\),
\(g\in
\underline{\operatorname{Aut}}_{\mathcal O_{S'}\text{-mod.}}(F_{S'})\),
and}
\[
 X,Y\in
 \underline{\operatorname{Lie}}\!
  \left(\underline{\operatorname{Aut}}_{\mathcal O_S\text{-mod.}}(F)/S\right)
  (S')
 =
 \underline{\operatorname{End}}_{\mathcal O_{S'}\text{-mod.}}(F_{S'}).
\]

\textit{Proof.}\footnote{Editor's note: the original has been expanded
and simplified in what follows.}
After base change, we may assume \(S'=S\). Put \(I=I_S\) and
\(\mathcal O_I=\mathcal O_S[t]\), with \(t^2=0\). Recall from \SGARefLink{sga3:ii:scholium:4-5-1}{4.5.1}
that the inclusion
\[
 i:\underline{\operatorname{End}}_{\mathcal O_S\text{-mod.}}(F)
 \hookrightarrow
 \underline{\operatorname{Aut}}_{\mathcal O_I\text{-mod.}}(F_I)
\]
sends \(Y\) to \(\mathrm{id}+tY\). By the definition of
\(\operatorname{Ad}(g)\) in 4.1.A,
\[
 \mathrm{id}+t\operatorname{Ad}(g)(Y)
 =g\circ(\mathrm{id}+tY)\circ g^{-1}
 =\mathrm{id}+t(g\circ Y\circ g^{-1}),
\]
and therefore
\(\operatorname{Ad}(g)(Y)=g\circ Y\circ g^{-1}\).

Let \(I'\) be a second copy of \(I_S\), and put
\(\mathcal O_{I'}=\mathcal O_S[t']\), with \((t')^2=0\). Apply \SGARefLink{sga3:ii:subsection:4-7-3}{4.7.3}
to
\[
 G=\underline{\operatorname{Aut}}_{\mathcal O_S\text{-mod.}}(F),
 \qquad
 L=\underline{\operatorname{Lie}}(G/S)
  =\underline{\operatorname{End}}_{\mathcal O_S\text{-mod.}}(F).
\]
Here \(L=\operatorname{End}(F)\) is the differential identification of
\SGARefLink{sga3:ii:proposition:4-5}{Proposition 4.5}. The re-edition prints \(\operatorname{Aut}(F)\) at
this point, but \(\operatorname{End}(F)\) is the type-correct object
used by the ensuing endomorphism calculation. Identify \(X\) with its image under the
canonical section \(L(S)\hookrightarrow L(I)\). Its image in
\(G(I\times I')\) is \(\mathrm{id}+t'X\), whose inverse is
\(\mathrm{id}-t'X\). Similarly, \(Y\) lifts to
\(\mathrm{id}+tY\), with inverse \(\mathrm{id}-tY\). Their
commutator is
\[
\begin{aligned}
 &(\mathrm{id}+t'X)\circ(\mathrm{id}+tY)
 \circ(\mathrm{id}-t'X)\circ(\mathrm{id}-tY)\\
 &\hspace{5em}
 =\mathrm{id}+tt'(X\circ Y-Y\circ X).
\end{aligned}
\]
This is the image in \(G(I\times I')\) of
\(Z=X\circ Y-Y\circ X\in L(S)\): first \(Z\) maps to
\(tZ\in L(I)\), and then to
\(\mathrm{id}+t'tZ\in G(I\times I')\). By \SGARefLink{sga3:ii:subsection:4-7-3}{4.7.3},
\([X,Y]=X\circ Y-Y\circ X\).

\medskip
\noindent\SGAInlineRefTarget{sga3:ii:corollary:4-8-1}\textbf{Corollary 4.8.1.}
\textit{Let \(G\) be a good \(S\)-group and let
\(x,y,z\in\underline{\operatorname{Lie}}(G/S)(S')\). Then}
\[
 [x,[y,z]]+[y,[z,x]]+[z,[x,y]]=0.
\]
\footnote{Editor's note: the original has been expanded in what follows.}
Indeed, since \(G\) is good,
\(\underline{\operatorname{Lie}}(G/S)\) is a good
\(\mathcal O_S\)-module; hence \SGARefLink{sga3:ii:theorem:4-7}{Theorem 4.7} shows that
\[
 \underline{\operatorname{Aut}}_{\mathcal O_S\text{-mod.}}
   \bigl(\underline{\operatorname{Lie}}(G/S)\bigr)
\]
is a good \(S\)-group. The morphism
\[
 \operatorname{Ad}:G\longrightarrow
 \underline{\operatorname{Aut}}_{\mathcal O_S\text{-mod.}}
  \bigl(\underline{\operatorname{Lie}}(G/S)\bigr)
\]
and the functoriality in \SGARefLink{sga3:ii:subsection:4-7-3:item:i}{4.7.3(i)} yield
\[
 \operatorname{ad}[x,y]
 =[\operatorname{ad}x,\operatorname{ad}y].
\]
Together with \SGARefLink{sga3:ii:proposition:4-8}{Proposition 4.8}, this gives
\[
 \operatorname{ad}[x,y]
 =\operatorname{ad}x\circ\operatorname{ad}y
  -\operatorname{ad}y\circ\operatorname{ad}x.
\]
Applying both sides to \(z\) gives the Jacobi identity.

\medskip
\noindent\SGAInlineRefTarget{sga3:ii:corollary:4-8-2}\textbf{Corollary 4.8.2.}
\textit{Let \(G\) be a good \(S\)-group acting linearly on a good
\(\mathcal O_S\)-module \(F\); equivalently, let \(F\) be a
\(G\)-\(\mathcal O_S\)-module, with \(G\) and \(F\) good. Then the
linear map}
\[
 \rho':\underline{\operatorname{Lie}}(G/S)
 \longrightarrow
 \underline{\operatorname{End}}_{\mathcal O_S\text{-mod.}}(F)
\]
\textit{is a representation; that is,}
\[
 \rho'([x,y])
 =\rho'(x)\circ\rho'(y)-\rho'(y)\circ\rho'(x).
\]

\medskip
\noindent\SGAInlineRefTarget{sga3:ii:scholium:4-9}\textbf{Scholium 4.9.}
To every good \(S\)-group \(G\) (for example, every representable
one) we have functorially associated a good \(\mathcal O_S\)-module
\(\underline{\operatorname{Lie}}(G/S)\) equipped with a bilinear map
satisfying
\[
 [x,y]+[y,x]=0,
 \qquad
 [x,[y,z]]+[y,[z,x]]+[z,[x,y]]=0.
\]
We shall call \(\underline{\operatorname{Lie}}(G/S)\), with this
structure, the \emph{Lie algebra} of \(G\) over \(S\). The qualification
is that it is not yet known whether it is an
\(\mathcal O_S\)-Lie algebra in the strict sense.
Every linear representation of \(G\) in a good
\(\mathcal O_S\)-module \(F\) induces a representation of its Lie
algebra. In particular, the adjoint representation of \(G\) induces
the adjoint representation of its Lie algebra.

\medskip
\noindent\SGAInlineRefTarget{sga3:ii:definition:4-10}\textbf{Definition 4.10.}
An \(S\)-functor in groups \(G\) is called \emph{very good} if it is
good and \(\underline{\operatorname{Lie}}(G/S)\) is an
\(\mathcal O_S\)-Lie algebra, that is, if \([x,x]=0\) identically.

\medskip
\noindent\SGAInlineRefTarget{sga3:ii:example:4-10-1}\textbf{Examples 4.10.1.}
The following \(S\)-groups are very good:
\(\underline{\operatorname{Aut}}_{\mathcal O_S\text{-mod.}}(F)\)
for every good \(\mathcal O_S\)-module \(F\) (by \SGARefLink{sga3:ii:theorem:4-7}{4.7} and \SGARefLink{sga3:ii:proposition:4-8}{4.8});
every representable group (as shown below); every good \(S\)-group
admitting a monomorphism into a very good \(S\)-group, in particular
every good subgroup functor of a representable group; and every good
\(S\)-group admitting a faithful linear representation in a good
\(\mathcal O_S\)-module, for example every good \(S\)-group for which
\(\operatorname{Ad}\) is a monomorphism.

\subsection*{\SGAInlineRefTarget{sga3:ii:subsection:4-11}4.11.}
Suppose now that \(G\) is a group scheme over \(S\). By 4.1.4,
\(\underline{\operatorname{Lie}}(G/S)(S)\) identifies with the group
of infinitesimal automorphisms of \(G/S\) invariant under right
translation; by \SGARefLink{sga3:ii:subsection:3-14}{3.14}, this is the group of derivations of
\(\mathcal O_G\) over \(\mathcal O_S\) invariant under right
translation. This identification respects the module structure and
is an anti-isomorphism of Lie algebras.

To see this, argue as in \SGARefLink{sga3:ii:subsection:4-7-3}{4.7.3}. Put
\[
 \mathcal O_I=\mathcal O_S[t],
 \qquad
 \mathcal O_{I'}=\mathcal O_S[t'],
 \qquad
 t^2=0=(t')^2,
\]
and let \(x\in L(I)\), \(y\in L(I')\), where
\(L=\underline{\operatorname{Lie}}(G/S)\). Left translation
\(\lambda_x\), respectively \(\lambda_y\), is an \(S\)-automorphism
of \(G_{I\times I'}\) inducing the identity on \(G_{I'}\),
respectively on \(G_I\). It corresponds to the
\(\mathcal O_S\)-automorphism
\[
 u=\mathrm{id}+t\,d_x,
 \qquad\text{respectively}\qquad
 v=\mathrm{id}+t'\,d_y
\]
of
\[
 \mathcal O_{G_{I\times I'}}
 =\mathcal O_G[t,t']/(t^2,(t')^2),
\]
where \(d_x,d_y\) are \(\mathcal O_S\)-derivations of
\(\mathcal O_G\) invariant under right translation. The
correspondence between \(S\)-automorphisms of \(G_{I\times I'}\)
and \(\mathcal O_S\)-automorphisms of
\(\mathcal O_{G_{I\times I'}}\) is contravariant. Therefore
\[
 \lambda_x\lambda_y\lambda_x^{-1}\lambda_y^{-1}
\]
corresponds to
\[
 v^{-1}u^{-1}vu
 =\mathrm{id}+tt'(d_yd_x-d_xd_y).
\]
It follows from \SGARefLink{sga3:ii:subsection:4-7-3}{4.7.3} that
\[
 x\longmapsto-d_x
\]
is an isomorphism of Lie algebras; for further details, see [DG70],
\(\S\)II.4, 4.4 and \SGARefLink{sga3:ii:definition:4-6}{4.6}. The same argument applies to
\[
 \underline{\operatorname{Lie}}(G/S)(S')
 =
 \underline{\operatorname{Lie}}(G_{S'}/S')(S')
\]
for every \(S'\to S\). Thus one recovers the classical
definition.\footnote{Editor's note: the original has been expanded in
what follows. In particular, the numbering \SGARefLink{sga3:ii:scholium:4-11-1}{4.11.1}--\SGARefLink{sga3:ii:corollary:4-11-8}{4.11.8} has been
added for later reference.}

\medskip
\noindent\SGAInlineRefTarget{sga3:ii:scholium:4-11-1}\textbf{Scholium 4.11.1.}
\textit{Via the isomorphism \(x\mapsto-d_x\),
\(\underline{\operatorname{Lie}}(G/S)\) identifies with the functor
that assigns to each \(S'\to S\) the \(\mathcal O(S')\)-Lie algebra
of derivations of \(G_{S'}\) relative to \(S'\) invariant under right
translation.}

Since every representable group is already known to be good by
\SGARefLink{sga3:ii:example:4-6-1}{4.6.1}, the preceding discussion gives:

\medskip
\noindent\SGAInlineRefTarget{sga3:ii:corollary:4-11-2}\textbf{Corollary 4.11.2.}
\textit{Every representable group is very good.}

Let \(\varepsilon:S\to G\) be the unit section of \(G\), and put
\[
 \omega^1_{G/S}=\varepsilon^*(\Omega^1_{G/S}).
\]
Recall from \SGARefLink{sga3:ii:proposition:3-3}{3.3} that
\(\underline{\operatorname{Lie}}(G/S)\) is represented by the vector
fibration \(V(\omega^1_{G/S})\).

\medskip
\noindent\SGAInlineRefTarget{sga3:ii:scholium:4-11-3}\textbf{Scholium 4.11.3.}
\textit{Thus to every \(S\)-group scheme \(G\) we have functorially
associated a vector fibration}
\[
 \operatorname{Lie}(G/S)=V(\omega^1_{G/S})
\]
\textit{over \(S\), representing the functor
\(\underline{\operatorname{Lie}}(G/S)\), and hence endowed with the
structure of an \(S\)-scheme in \(\mathcal O_S\)-Lie algebras.
Moreover, this construction commutes with base extension and finite
products (cf. \SGARefLink{sga3:ii:proposition:3-4}{3.4} and \SGARefLink{sga3:ii:proposition:3-8}{3.8}).}

\medskip
\noindent\SGAInlineRefTarget{sga3:ii:remark:4-11-4}\textbf{Remarks 4.11.4.}
\footnote{Editor's note: what follows has been expanded, taking account
of the additions made in Expose I, \(\S\)\SGARefLink{sga3:i:subsection:6-8}{6.8}.}
Let \(\pi:G\to S\) be the structural morphism.

\textit{(a)} The \(\mathcal O_G\)-module \(\Omega^1_{G/S}\) is
naturally \((G\times_SG)\)-equivariant (cf. I, \(\S\)6). Hence, by
\SGARefLink{sga3:i:proposition:6-8-1}{I~6.8.1},
\[
 \Omega^1_{G/S}\simeq\pi^*(\omega^1_{G/S}).
\]
For example, \(\Omega^1_{G/S}\) is locally free, respectively locally
free of finite type, whenever \(\omega^1_{G/S}\) is. This holds in
particular when \(S\) is the spectrum of a field, respectively when
\(S\) is the spectrum of a field and \(G\) is locally of finite type
over \(S\).

\textit{(b)} By \SGARefLink{sga3:i:remark:6-8-2}{I~6.8.2}, \(\omega^1_{G/S}\) has a canonical
\(G\)-\(\mathcal O_S\)-module structure, which induces the adjoint
action on
\[
 V(\omega^1_{G/S})=\operatorname{Lie}(G/S).
\]
\footnote{Editor's note: for this reason, the linear action of \(G\) on
\(\omega^1_{G/S}\) could be called the pre-adjoint action; by an abuse
of terminology it will still be called the adjoint action. A slightly
more general construction, used in Expose III (especially \SGARefLink{sga3:iii:paragraph:0-8}{III~0.8}), is
as follows. Suppose that to every \(Y\to S\) one assigns an
\(\mathcal O_Y\)-module \(\mathcal N(Y)\), and to every
\(S\)-morphism \(\phi:Z\to Y\) a functorial morphism of
\(\mathcal O_Z\)-modules
\(\mathcal N(\phi):\phi^*\mathcal N(Y)\to\mathcal N(Z)\).
For example, if \(\mathcal J\) is a quasi-coherent ideal of
\(\mathcal O_S\), one may take
\(\mathcal N(Y)=\mathcal J\mathcal O_Y\). The action of \(G\) on
\(\omega^1_{G/S}\) then induces a generalized adjoint action of \(G\)
on the \(S\)-functor in abelian groups that assigns to
\(f:Y\to S\) the group
\(\operatorname{Hom}_{\mathcal O_Y}
(f^*\omega^1_{G/S},\mathcal N(Y))\).}

\textit{(c)} By EGA~I, 5.3.11 and 5.4.6,
\(\varepsilon:S\to G\) is an immersion; it is a closed immersion if
\(G\) is separated over \(S\). Thus \(\omega^1_{G/S}\) identifies
with \(\mathcal I/\mathcal I^2\), where \(\mathcal I\) is the
quasi-coherent ideal defining \(\varepsilon(S)\) in an open
subscheme \(U\subset G\) in which \(\varepsilon(S)\) is closed. If
\(G\to S\) is separated, one may take \(U=G\). If
\(G=\operatorname{Spec}A(G)\) is affine over \(S\), then
\(\mathcal I\) is the augmentation ideal of \(A(G)\), namely the
kernel of
\[
 \varepsilon^\natural:A(G)\longrightarrow\mathcal O_S
\]
(cf. \SGARefLink{sga3:i:subsection:4-2}{I~4.2}).\footnote{Editor's note: the preceding description of
\(\omega^1_{G/S}\), which is useful later (for example in
VII\textsubscript{A}, \SGARefLink{sga3:VIIA:5.5}{5.5}), has been added.}

\medskip
\noindent\SGAInlineRefTarget{sga3:ii:remark:4-11-5}\textbf{Remark 4.11.5.}
The isomorphism
\(\Omega^1_{G/S}\simeq\pi^*(\omega^1_{G/S})\) identifies the
\(\mathcal O_S\)-module \(\omega^1_{G/S}\) with the sheaf
\[
 \pi_*(\Omega^1_{G/S})^G
\]
of differentials of \(G\) relative to \(S\) invariant under right
translation: its sections over an open \(U\subset S\) are the
sections of \(\Omega^1_{G/S}\) over \(\pi^{-1}(U)\) invariant under
right translation (cf. \SGARefLink{sga3:i:remark:6-8-3}{I~6.8.3}; compare also
VII\textsubscript{A}, \SGARefLink{sga3:VIIA:2.4}{2.4}).

\medskip
\noindent\SGAInlineRefTarget{sga3:ii:notation:4-11-6}\textbf{Notation 4.11.6.}
Write \(\mathop{\mathscr L\!ie}(G/S)\) for the sheaf of sections of the vector
fibration \(\operatorname{Lie}(G/S)\to S\). It is the
\(\mathcal O_S\)-module
\[
 \mathop{\mathscr L\!ie}(G/S)
 =(\omega^1_{G/S})^\vee
 =\mathcal Hom_{\mathcal O_S}(\omega^1_{G/S},\mathcal O_S),
\]
dual to \(\omega^1_{G/S}\) (cf. EGA~II, 1.7.9), and it carries a
structure of \(\mathcal O_S\)-Lie algebra.

In general this construction does not commute with base extension.
Consequently, the Lie-algebra structure on this module does not
determine the structure of an \(S\)-scheme in
\(\mathcal O_S\)-Lie algebras on
\(\operatorname{Lie}(G/S)\).

Nevertheless:

\medskip
\noindent\SGAInlineRefTarget{sga3:ii:lemma:4-11-7}\textbf{Lemma 4.11.7.}
\textit{Suppose that \(\omega^1_{G/S}\) is locally free of finite type.
This is the case, in particular, if \(G\) is smooth over \(S\)
(cf. EGA~IV\textsubscript{4}, 17.2.4), or if \(S\) is the spectrum
of a field and \(G\) is locally of finite type over \(S\). Then}
\[
 \mathop{\mathscr L\!ie}(G/S)^\vee
 \simeq(\omega^1_{G/S})^{\vee\vee}
 \simeq\omega^1_{G/S},
\]
\textit{and therefore}
\[
 \operatorname{Lie}(G/S)
 =V(\omega^1_{G/S})
 =V\bigl(\mathop{\mathscr L\!ie}(G/S)^\vee\bigr)
 =W\bigl(\mathop{\mathscr L\!ie}(G/S)\bigr),
\]
\textit{where the final equality follows from \SGARefLink{sga3:i:corollary:4-6-5-1}{I~4.6.5.1}.}

Finally, let \(G\to H\) be a monomorphism of functors in groups.
Then
\[
 \operatorname{Lie}(G/S)\longrightarrow\operatorname{Lie}(H/S)
\]
is also a monomorphism (cf. \SGARefLink{sga3:ii:proposition:3-7}{3.7}). Every vector monomorphism of vector
fibrations is a closed immersion.
It follows that:

\medskip
\noindent\SGAInlineRefTarget{sga3:ii:corollary:4-11-8}\textbf{Corollary 4.11.8.}
\textit{Let \(G\hookrightarrow H\) be a monomorphism of
\(S\)-group schemes.}

\textit{(i) The morphism}
\[
 \operatorname{Lie}(G/S)\longrightarrow\operatorname{Lie}(H/S)
\]
\textit{is a closed immersion; consequently}
\[
 \omega^1_{H/S}\longrightarrow\omega^1_{G/S}
\]
\textit{is an epimorphism.}

\textit{(ii) If \(\omega^1_{G/S}\) is locally free of finite type,
then the corresponding morphism}
\[
 \mathop{\mathscr L\!ie}(G/S)\longrightarrow
 \mathop{\mathscr L\!ie}(H/S)
\]
\textit{identifies \(\mathop{\mathscr L\!ie}(G/S)\) with an
\(\mathcal O_S\)-submodule of \(\mathop{\mathscr L\!ie}(H/S)\) that is locally
a direct summand.}

